\begin{exbox}[两侧虚通道的共同位移与差分位移]
在 $0\leftrightarrow1$ 共振处，保留最近的 $|-1\rangle,|2\rangle$，
并取常数实耦合 $g_B$（单位角频率），令 $\Delta_Q=2\omega_{\rm r}$。
按 $|-1\rangle,|0\rangle,|1\rangle,|2\rangle$ 排列，四态模型是
\[
\frac H\hbar=\begin{pmatrix}
\Delta_Q&g_B&0&0\\g_B&0&g_B&0\\
0&g_B&0&g_B\\0&0&g_B&\Delta_Q
\end{pmatrix}.
\]
求能量依赖的两态自能及低能劈裂。

\textbf{解。}取 $P$ 保留中间两态，$Q$ 保留外部两态。以 $\varepsilon=z/\hbar$ 表示能量对应的角频率，
分块消元立即给
\begin{equation}
\frac{H_{\rm eff}(z)}{\hbar}
=\begin{pmatrix}0&g_B\\g_B&0\end{pmatrix}
+\frac{g_B^2}{\varepsilon-\Delta_Q}I_2.
\label{eq:feqo-four-state-schur-effective}
\end{equation}
这里分母是角频率；如果使用能量变量 $z$，分子必须是 $\hbar g_B^2$，
不能把 $g_B^2/(z-\hbar\Delta_Q)$ 直接称为频率。
在 $|g_B|\ll\Delta_Q$、$|\varepsilon|\ll\Delta_Q$ 时，
\[
\frac{H_{\rm eff}}\hbar
=g_B\sigma_x-\frac{g_B^2}{\Delta_Q}I_2
+O(g_B^3/\Delta_Q^2).
\]
二阶位移完全相同，故只产生公共相位；它不移动这两态的相对共振。
若两侧耦合为 $g_-,g_+$，则二阶位移分别为
$-|g_-|^2/\Delta_Q$ 与 $-|g_+|^2/\Delta_Q$，差值才改变共振失谐。

对对称模型还可核对第三阶。保留块的对称/反对称本征态对应 $\eta=\pm1$，
低能极点满足 $\varepsilon_\eta=\eta g_B+g_B^2/(\varepsilon_\eta-\Delta_Q)$，故
\[
\varepsilon_\eta=\frac{\Delta_Q+\eta g_B-
\sqrt{(\Delta_Q-\eta g_B)^2+4g_B^2}}2
=\eta g_B-\frac{g_B^2}{\Delta_Q}
-\eta\frac{g_B^3}{\Delta_Q^2}+O(g_B^4/\Delta_Q^3).
\]
所以有效振荡频率为
$\varepsilon_+-\varepsilon_-=2g_B[1-g_B^2/\Delta_Q^2+O((g_B/\Delta_Q)^4)]$。
外部本征振幅相对目标幅度是 $g_B/(\varepsilon_\eta-\Delta_Q)$，
说明虚通道人口为 $O((g_B/\Delta_Q)^2)$。
过长脉冲还会积累第三阶频率修正，不能只检查一瞬间的泄漏。
\end{exbox}

\begin{exbox}[晶体两束衍射的厚度标尺]
设某一反射在精确两束共振处，势 Fourier 系数的模长为给定参数 $|V_G|$。
固定速度 $v_0$ 时，以厚度 $z$ 为演化变量：
\[
\ii\hbar v_0\partial_z\binom{a_0}{a_G}
=\begin{pmatrix}0&V_G^*\\V_G&0\end{pmatrix}\binom{a_0}{a_G}.
\]
求从零束入射后的转移概率和特征厚度。

\textbf{解。}直接指数化给
\begin{equation}
P_G(z)=\sin^2\left(\frac{|V_G|z}{\hbar v_0}\right),
\quad\Delta E_{\rm gap}=2|V_G|,
\quad\xi_G=\frac{\pi\hbar v_0}{|V_G|}.
\label{eq:feqo-two-beam-bandgap}
\end{equation}
给定 $K_0=\SI{200}{\kilo\electronvolt}$、$|V_G|=\SI{5.9}{\electronvolt}$，
$\hbar v_0=1.37205\times10^{-7}\,\si{\electronvolt\metre}$，
所以 $\xi_G=\SI{73.06}{\nano\metre}$。
首次半转移厚度为 $\xi_G/4=\SI{18.27}{\nano\metre}$，
首次全转移厚度为 $\xi_G/2=\SI{36.53}{\nano\metre}$。
此计算使用给定势系数；其它反射被激发或多束谱隙不够时，应返回完整 Bloch 纤维。
\end{exbox}
\begin{figure}[htbp]
 \centering
 \includegraphics[width=0.76\linewidth]{pinem_recoil_breaking.pdf}
 \caption{反冲可分辨时 Bessel 谱失效。作用时间归一化为 $0\le\tau\le1$，
 无量纲链为 $\ii\partial_\tau c_\ell=\chi\ell^2c_\ell+2.5(c_{\ell-1}+c_{\ell+1})$。
 曲线比较 $\chi=0,0.20,0.40$，初态为中心格点。
 数值截断到 $|\ell|\le12$；这里展示谱形变化，完整截断误差控制见前面的 Duhamel 界及收敛例题。}
 \label{fig:pinem-recoil-breaking}
\end{figure}
